\chapter{Chuẩn hóa các lược đồ quan hệ về dạng chuẩn 3NF}

\section{Phân tích các lược đồ quan hệ về dạng chuẩn 3NF}
Phân tích dưới đây áp dụng cho 20 lược đồ đã trình bày ở Chương 3. Với mỗi bảng, tập phụ thuộc hàm được viết theo khóa chính đã chọn trong mô hình. Các thuộc tính được xem là giá trị đơn theo mô tả lược đồ nên đạt 1NF. Với khóa đơn, không phát sinh phụ thuộc bộ phận; với khóa ghép, các thuộc tính không khóa được gắn với toàn bộ khóa. Không có phụ thuộc bắc cầu nào được nêu trong tập phụ thuộc của mô hình, vì vậy các lược đồ đạt 3NF theo các phụ thuộc được xác định ở đây.

Đây là phân tích dựa trên lược đồ và khóa của báo cáo. Muốn chứng minh đầy đủ các phụ thuộc hàm của hệ thống thực tế cần xác nhận thêm các quy tắc nghiệp vụ và ràng buộc duy nhất đang áp dụng.

\begin{description}[leftmargin=4.1cm,labelwidth=3.9cm,style=nextline]
  \item[1. \texttt{stores}] Lược đồ: \texttt{(id, code, name, address, status)}. \textbf{F:} \texttt{id} $\rightarrow$ \texttt{code, name, address, status}. \textbf{Khóa:} \texttt{id}; \texttt{code} là khóa duy nhất thay thế. Tất cả thuộc tính không khóa phụ thuộc vào khóa; không xác định phụ thuộc bắc cầu. Lược đồ đạt 1NF, 2NF và 3NF.
  \item[2. \texttt{employees}] Lược đồ: \texttt{(id, username, full\_name, phone, role, status)}. \textbf{F:} \texttt{id} $\rightarrow$ \texttt{username, full\_name, phone, role, status}. \textbf{Khóa:} \texttt{id}; \texttt{username} là khóa duy nhất thay thế. Lược đồ đạt 1NF, 2NF và 3NF theo tập phụ thuộc đã nêu.
  \item[3. \texttt{employee\_store\_assignments}] Lược đồ: \texttt{(employee\_id, store\_id, granted\_by, granted\_at)}. \textbf{F:} \texttt{(employee\_id, store\_id)} $\rightarrow$ \texttt{granted\_by, granted\_at}. \textbf{Khóa:} \texttt{(employee\_id, store\_id)}. Các thuộc tính phân công phụ thuộc vào cả hai thành phần khóa; lược đồ đạt 1NF, 2NF và 3NF.
  \item[4. \texttt{customers}] Lược đồ: \texttt{(id, code, full\_name, phone, email, status)}. \textbf{F:} \texttt{id} $\rightarrow$ \texttt{code, full\_name, phone, email, status}. \textbf{Khóa:} \texttt{id}; \texttt{code} và \texttt{phone} được khai báo duy nhất trong mô hình. Lược đồ đạt 1NF, 2NF và 3NF theo các phụ thuộc này.
  \item[5. \texttt{accounts}] Lược đồ: \texttt{(id, customer\_id, code, username, balance, status)}. \textbf{F:} \texttt{id} $\rightarrow$ \texttt{customer\_id, code, username, balance, status}. \textbf{Khóa:} \texttt{id}; \texttt{customer\_id}, \texttt{code} và \texttt{username} là các khóa duy nhất thay thế theo lược đồ. Lược đồ đạt 1NF, 2NF và 3NF.
  \item[6. \texttt{machines}] Lược đồ: \texttt{(id, store\_id, code, room\_type, status, hourly\_rate)}. \textbf{F:} \texttt{id} $\rightarrow$ \texttt{store\_id, code, room\_type, status, hourly\_rate}. \textbf{Khóa:} \texttt{id}; \texttt{(store\_id, code)} là khóa duy nhất thay thế. Lược đồ đạt 1NF, 2NF và 3NF.
  \item[7. \texttt{usage\_sessions}] Lược đồ: \texttt{(id, customer\_id, machine\_id, opened\_by, closed\_by, started\_at, ended\_at, cost, status)}. \textbf{F:} \texttt{id} $\rightarrow$ \texttt{customer\_id, machine\_id, opened\_by, closed\_by, started\_at, ended\_at, cost, status}. \textbf{Khóa:} \texttt{id}. Không có phụ thuộc bộ phận hoặc bắc cầu được xác định; lược đồ đạt 1NF, 2NF và 3NF.
  \item[8. \texttt{account\_transactions}] Lược đồ: \texttt{(id, account\_id, employee\_id, session\_id, store\_id, promotion\_id, type, amount, created\_at)}. \textbf{F:} \texttt{id} $\rightarrow$ \texttt{account\_id, employee\_id, session\_id, store\_id, promotion\_id, type, amount, created\_at}. \textbf{Khóa:} \texttt{id}. Lược đồ đạt 1NF, 2NF và 3NF theo tập phụ thuộc đã nêu.
  \item[9. \texttt{invoices}] Lược đồ: \texttt{(id, session\_id, customer\_id, employee\_id, store\_id, usage\_cost, total\_amount, status, issued\_at)}. \textbf{F:} \texttt{id} $\rightarrow$ \texttt{session\_id, customer\_id, employee\_id, store\_id, usage\_cost, total\_amount, status, issued\_at}. \textbf{Khóa:} \texttt{id}; \texttt{session\_id} là duy nhất nên là khóa thay thế. Lược đồ đạt 1NF, 2NF và 3NF.
  \item[10. \texttt{promotions}] Lược đồ: \texttt{(id, store\_id, code, minimum\_amount, bonus\_percent, starts\_at, ends\_at, status)}. \textbf{F:} \texttt{id} $\rightarrow$ \texttt{store\_id, code, minimum\_amount, bonus\_percent, starts\_at, ends\_at, status}. \textbf{Khóa:} \texttt{id}; \texttt{(store\_id, code)} là khóa duy nhất thay thế. Lược đồ đạt 1NF, 2NF và 3NF.
  \item[11. \texttt{management\_audit}] Lược đồ: \texttt{(id, employee\_id, store\_id, action, entity\_type, entity\_id, created\_at)}. \textbf{F:} \texttt{id} $\rightarrow$ \texttt{employee\_id, store\_id, action, entity\_type, entity\_id, created\_at}. \textbf{Khóa:} \texttt{id}. Lược đồ đạt 1NF, 2NF và 3NF theo phụ thuộc khóa chính.
  \item[12. \texttt{menu\_categories}] Lược đồ: \texttt{(id, store\_id, name, page, status)}. \textbf{F:} \texttt{id} $\rightarrow$ \texttt{store\_id, name, page, status}. \textbf{Khóa:} \texttt{id}; \texttt{(store\_id, name)} là khóa duy nhất thay thế. Lược đồ đạt 1NF, 2NF và 3NF.
  \item[13. \texttt{menu\_items}] Lược đồ: \texttt{(id, store\_id, category\_id, name, price, status)}. \textbf{F:} \texttt{id} $\rightarrow$ \texttt{store\_id, category\_id, name, price, status}. \textbf{Khóa:} \texttt{id}. Lược đồ đạt 1NF, 2NF và 3NF theo phụ thuộc khóa chính.
  \item[14. \texttt{menu\_orders}] Lược đồ: \texttt{(id, store\_id, account\_id, session\_id, machine\_id, handled\_by, status, total\_amount, created\_at)}. \textbf{F:} \texttt{id} $\rightarrow$ \texttt{store\_id, account\_id, session\_id, machine\_id, handled\_by, status, total\_amount, created\_at}. \textbf{Khóa:} \texttt{id}. Lược đồ đạt 1NF, 2NF và 3NF theo tập phụ thuộc đã nêu.
  \item[15. \texttt{menu\_order\_items}] Lược đồ: \texttt{(order\_id, item\_id, item\_name, unit\_price, quantity)}. \textbf{F:} \texttt{(order\_id, item\_id)} $\rightarrow$ \texttt{item\_name, unit\_price, quantity}. \textbf{Khóa:} \texttt{(order\_id, item\_id)}. Tên món và đơn giá được lưu như dữ liệu tại thời điểm đặt; các thuộc tính dòng đơn phụ thuộc vào toàn bộ khóa ghép. Lược đồ đạt 1NF, 2NF và 3NF theo phụ thuộc này.
  \item[16. \texttt{pos\_order\_reversals}] Lược đồ: \texttt{(id, order\_id, store\_id, requested\_by, decided\_by, status, reason, created\_at)}. \textbf{F:} \texttt{id} $\rightarrow$ \texttt{order\_id, store\_id, requested\_by, decided\_by, status, reason, created\_at}. \textbf{Khóa:} \texttt{id}; \texttt{order\_id} là duy nhất thay thế. Lược đồ đạt 1NF, 2NF và 3NF.
  \item[17. \texttt{menu\_promotions}] Lược đồ: \texttt{(id, store\_id, item\_id, code, kind, value, starts\_at, ends\_at)}. \textbf{F:} \texttt{id} $\rightarrow$ \texttt{store\_id, item\_id, code, kind, value, starts\_at, ends\_at}. \textbf{Khóa:} \texttt{id}. Lược đồ đạt 1NF, 2NF và 3NF theo phụ thuộc khóa chính.
  \item[18. \texttt{menu\_order\_pricing}] Lược đồ: \texttt{(order\_id, promotion\_id, discount\_amount, final\_amount)}. \textbf{F:} \texttt{order\_id} $\rightarrow$ \texttt{promotion\_id, discount\_amount, final\_amount}. \textbf{Khóa:} \texttt{order\_id}. Lược đồ đạt 1NF, 2NF và 3NF theo phụ thuộc khóa chính.
  \item[19. \texttt{customer\_notification\_events}] Lược đồ: \texttt{(id, store\_id, title, message, event\_type, created\_at)}. \textbf{F:} \texttt{id} $\rightarrow$ \texttt{store\_id, title, message, event\_type, created\_at}. \textbf{Khóa:} \texttt{id}. Lược đồ đạt 1NF, 2NF và 3NF theo phụ thuộc khóa chính.
  \item[20. \texttt{customer\_notification\_reads}] Lược đồ: \texttt{(event\_id, customer\_id, read\_at)}. \textbf{F:} \texttt{(event\_id, customer\_id)} $\rightarrow$ \texttt{read\_at}. \textbf{Khóa:} \texttt{(event\_id, customer\_id)}. Thời điểm đọc phụ thuộc vào toàn bộ khóa ghép; lược đồ đạt 1NF, 2NF và 3NF.
\end{description}

\section{Các lược đồ sau khi chuẩn hóa}
Dựa trên các khóa và phụ thuộc hàm đã nêu, không cần tách thêm thuộc tính khỏi 20 bảng của mô hình báo cáo để đạt 3NF. Các khóa chính và khóa ngoại được giữ như Chương 3; các lược đồ sau chuẩn hóa vì vậy vẫn gồm:

\begin{description}[leftmargin=4.1cm,labelwidth=3.9cm,style=nextline]
  \item[Nhóm cửa hàng và nhân viên] \texttt{stores(id, code, name, address, status)}; \texttt{employees(id, username, full\_name, phone, role, status)}; \texttt{employee\_store\_assignments(employee\_id, store\_id, granted\_by, granted\_at)}; \texttt{management\_audit(id, employee\_id, store\_id, action, entity\_type, entity\_id, created\_at)}.
  \item[Nhóm khách hàng và phiên máy] \texttt{customers(id, code, full\_name, phone, email, status)}; \texttt{accounts(id, customer\_id, code, username, balance, status)}; \texttt{machines(id, store\_id, code, room\_type, status, hourly\_rate)}; \texttt{usage\_sessions(id, customer\_id, machine\_id, opened\_by, closed\_by, started\_at, ended\_at, cost, status)}; \texttt{account\_transactions(id, account\_id, employee\_id, session\_id, store\_id, promotion\_id, type, amount, created\_at)}; \texttt{invoices(id, session\_id, customer\_id, employee\_id, store\_id, usage\_cost, total\_amount, status, issued\_at)}; \texttt{promotions(id, store\_id, code, minimum\_amount, bonus\_percent, starts\_at, ends\_at, status)}.
  \item[Nhóm thực đơn và đơn hàng] \texttt{menu\_categories(id, store\_id, name, page, status)}; \texttt{menu\_items(id, store\_id, category\_id, name, price, status)}; \texttt{menu\_orders(id, store\_id, account\_id, session\_id, machine\_id, handled\_by, status, total\_amount, created\_at)}; \texttt{menu\_order\_items(order\_id, item\_id, item\_name, unit\_price, quantity)}; \texttt{pos\_order\_reversals(id, order\_id, store\_id, requested\_by, decided\_by, status, reason, created\_at)}; \texttt{menu\_promotions(id, store\_id, item\_id, code, kind, value, starts\_at, ends\_at)}; \texttt{menu\_order\_pricing(order\_id, promotion\_id, discount\_amount, final\_amount)}.
  \item[Nhóm thông báo] \texttt{customer\_notification\_events(id, store\_id, title, message, event\_type, created\_at)}; \texttt{customer\_notification\_reads(event\_id, customer\_id, read\_at)}.
\end{description}

\section{Kết luận chương}
Theo cấu trúc thuộc tính, các khóa và tập phụ thuộc hàm đã nêu, 20 lược đồ quan hệ của báo cáo đạt 3NF và được giữ nguyên sau phân tích. Các bảng liên kết dùng khóa ghép; tên món và đơn giá trong dòng đơn được lưu tại thời điểm đặt để giữ lịch sử giao dịch. Kết luận về các phụ thuộc hàm còn cần được đối chiếu với quy tắc nghiệp vụ được xác nhận khi triển khai thực tế.
